% paper-macros.tex -- preamble snippets for an algorithm engineering paper.
% \input this after the document class. Remove what your class already
% provides (LIPIcs, for example, already defines theorem environments).

\usepackage{xcolor}
\usepackage{xspace}
\usepackage{graphicx}
\usepackage{booktabs}          % \toprule, \midrule, \bottomrule
\usepackage{siunitx}           % aligned numbers, units, thin-space grouping
\usepackage{amsmath,amssymb}
\usepackage{amsthm}            % skip if the class defines theorems
\usepackage{algorithm2e}       % or algorithmicx; pick one for the paper
\usepackage[capitalise,noabbrev]{cleveref} % load last

% If you use the paper-figures skill, load its palette here so tables and
% text use the same algorithm colors (adjust the path to your layout):
% \input{figures/palette/colors}
\providecolor{figWarn}{HTML}{D55E00} % color of \TODO and \GAP markers

% --- Algorithm and problem names ---------------------------------------------
% One macro per name. The macro fixes the spelling and capitalization once,
% so "RAPTOR" never becomes "Raptor" on page 7.
\newcommand{\algname}[1]{\textsc{#1}}
\newcommand{\Ours}{\algname{Ours}\xspace} % rename to your algorithm
% Example: \newcommand{\Raptor}{RAPTOR\xspace}

% --- Numbers -----------------------------------------------------------------
\sisetup{
  group-separator={\,},        % 10\,000 and 5\,132\,672
  group-minimum-digits=4,
  detect-weight=true,
  round-mode=none,             % never round silently; round in the data
}
% Usage: \num{10000} queries, \qty{7.3}{\milli\second},
% table columns of type S[table-format=4.1].

% --- Theorem environments (skip if the class defines them) --------------------
\theoremstyle{plain}
\newtheorem{theorem}{Theorem}
\newtheorem{lemma}[theorem]{Lemma}
\newtheorem{corollary}[theorem]{Corollary}
\newtheorem{proposition}[theorem]{Proposition}
\theoremstyle{definition}
\newtheorem{definition}[theorem]{Definition}
\newtheorem{observation}[theorem]{Observation}
\theoremstyle{remark}
\newtheorem*{remark}{Remark}

% --- Placeholders --------------------------------------------------------------
% Named \TODO, not \todo, to avoid a clash with the todonotes package.
% Visible marker for anything that still needs a fact, a number, or a proof
% step. Search for "TODO" before submission; the PDF must not contain any.
\newcommand{\TODO}[1]{\textcolor{figWarn}{[TODO: #1]}}
\newcommand{\GAP}[1]{\textcolor{figWarn}{[PROOF GAP: #1]}}

% --- Including figures ---------------------------------------------------------
% Every figure is built as its own standalone PDF (see assets/figures).
% Include the same PDF you previewed:
%   \begin{figure}[t]
%     \centering
%     \includegraphics{figures/build/scaling.pdf}
%     \caption{...}
%     \label{fig:scaling}
%   \end{figure}
% Do not scale figures with width=\linewidth; set the size in the figure
% source so that the font size in the figure matches the text.
